%%%%%%%%%%%%%%%%%%%%%%%%%%%%%%%%%%%%%%%%%%%%%%%%%%%%%%%%%%%%%%
%%%%%%%%%%%%%%%% MA-XXXX Análisis complejo %%%%%%%%%%%%%%%%%%%
%%%%%%%%%%%%%%%%%%%%%%%%%%%%%%%%%%%%%%%%%%%%%%%%%%%%%%%%%%%%%%
\newpage
\subsection*{MA-0702 Análisis Complejo}
\addcontentsline{toc}{subsection}{MA-0702 Análisis Complejo}

\hypertarget{MA0702}{}
\begin{refsection}
\begin{table}[H]
\centering{
\begin{tabular}{|ll|}
\hline
%\multicolumn{2}{|c|}{\textbf{Nivel de virtualidad del curso:} Virtual}\\ \hline
\textbf{Año:} IV  & \textbf{Requisitos:} Por determinar \\ 
\textbf{Ciclo:} VII  & \textbf{Correquisitos:} Ninguno \\ 
\textbf{Tipo de curso:} Teórico & \textbf{Horas presenciales por semana:}   \\ 
\textbf{Créditos:} 4 & \textbf{Horas de trabajo independiente por semana:}  \\ \hline
\end{tabular}
}
\end{table}



\subsubsection*{Descripción del curso}

Este es un curso del tercer año de la carrera Bach. y Lic. en Matemática en el que se presentan las bases del Análisis Complejo, que es una rama de las matemáticas que se enfoca en el estudio de funciones de una variable compleja. Parte de su importancia radica en su capacidad para extender y enriquecer nuestro entendimiento de las funciones, permitiendo el análisis y manipulación de fenómenos matemáticos y científicos más allá de lo que es posible con el análisis real. 

A través del análisis complejo se exploran propiedades únicas de los números complejos y de las funciones de variable compleja, en particular, se estudia cómo la existencia de derivadas complejas tiene consecuencias impresionantes para estas funciones, las cuales no son posibles en el contexto del análisis real.

Además, las herramientas que ofrece la teoría de funciones de variable compleja resultan tener tener aplicaciones fundamentales en prácticamente todas las ramas de la matemática moderna, así como a problemas de física, ingeniería y muchas otras disciplinas. También, las demostraciones de muchos de los teoremas más importantes de la matemática dependen de la teoría de funciones de variable compleja, como por ejemplo el Teorema Fundamental del Álgebra, el Teorema de los Números Primos y el Último Teorema de Fermat, por citar algunos. 



\subsubsection*{Objetivos}

\textbf{General:} Aplicar los conceptos y resultados básicos del análisis complejo para la resolución de problemas teóricos y prácticos. \\

\textbf{Específicos:} Al finalizar el curso se espera que la persona estudiante sea capaz de:

\begin{enumerate}
\item Comparar los conceptos, propiedades e implicaciones de la convergencia, continuidad, diferenciación e integración de funciones de variable real y funciones de variable compleja.
\item Demostrar resultados sobre límites, continuidad, convergencia, diferenciación e integración en el análisis de funciones de una variable compleja.
\item Aplicar los teoremas de diferenciación e integración de funciones de variable compleja en la demostración de resultados y la resolución de problemas teóricos y prácticos.
\item Relacionar aspectos algebraicos, analíticos y geométricos de los números complejos y de las funciones de variable compleja.
\item Indagar en la literatura sobre temas que permitan complementar su formación.
\item Comunicar ideas matemáticas de manera clara y efectiva en forma oral y escrita.
\end{enumerate}

\subsubsection*{Contenidos}

\begin{enumerate}
\item \textbf{Números complejos y funciones en el plano complejo (2.5 semanas):} 
\begin{enumerate}
    \item El cuerpo $\C$ de los números complejos, conjugados complejos y módulos. Forma polar de un número complejo. Representación de regiones en el plano complejo: rectas, semiplanos, regiones anulares.

    \item La esfera de Riemann $\C_{\infty}$: Proyección estereográfica y preservación de líneas y círculos.

    \item Topología de $\C$: sucesiones, límites, abiertos y cerrados, compacidad, conexidad y arco-conexidad. 

    \item Funciones en el plano complejo: continuidad, convergencia uniforme.
    
    \item Series de potencias, sus discos de convergencia.
    
    \item Ejemplos de funciones analíticas: exponenciales y  trigonométricas.
\end{enumerate}

\item \textbf{Diferenciación en el plano complejo (2 semanas):} 
\begin{enumerate}
    \item La derivada de una función compleja. Ejemplos.
    \item Ecuaciones de Cauchy-Riemann. 
    \item Funciones armónicas.
    \item Conformalidad de las funciones diferenciables.
    \item Aplicaciones conformes: transformaciones de Möbius.
\end{enumerate}

\item \textbf{Integración compleja, el Teorema de Cauchy y sus aplicaciones (4 semanas):} 
\begin{enumerate}
    \item Integrales de línea en $\C$.
    \item El Teorema de Goursat.
    \item La existencia local de primitivas y el Teorema de Cauchy en un disco.
    \item La fórmula integral de Cauchy.
    \item Las desigualdades de Cauchy, el Teorema de Liouville y sus consecuencias (el Teorema Fundamental del Algebra).
    \item Funciones holomorfas y representación local como series de potencias (analiticidad).
    \item El índice de una curva cerrada.
    \item La versión homotópica del Teorema de Cauchy y conexidad simple.
    \item Aplicaciones:
    \begin{itemize}
        \item El Teorema de Morera.
        \item Sucesiones de funciones holomorfas.
        \item Funciones holomorfas definidas en términos de integrales (ver e.g. \cite[\S 5.3]{Con78}).
        \item Continuación analítica y el principio de reflexión de Schwarz.
        \item El Teorema de aproximación de Runge.
    \end{itemize}
\end{enumerate}

\item \textbf{Singularidades, funciones meromorfas y el logaritmo complejo (3 semanas):} 
\begin{enumerate}
    \item Ceros, polos y clasificación de singularidades (polos, removibles y esenciales).
    \item El Teorema de Riemann  sobre singularidades removibles, el Teorema de Casorati-Weierstrass
    \item Series de Laurent.
    \item La fórmula del residuo y sus aplicaciones al cálculo de integrales y de series.
    \item El logaritmo complejo, ramas y otras funciones 
    \item El Principio del Argumento y aplicaciones: el teorema de Rouché, el teorema de la aplicación abierta, el principio del módulo máximo y el Lema de Schwarz.
\end{enumerate}

\item \textbf{Funciones enteras (1.5 semanas):} 
\begin{enumerate}
    \item La fórmula de Jensen.
    \item El orden de una función entera.
    \item Productos infinitos.
    \item Productos infinitos de Weierstrass. El ejemplo de $\sen{z}$.
\end{enumerate}

%\item \textbf{Las funciones Gama de Euler y Zeta de Riemann (X semanas):} 
%\begin{enumerate}
%    \item La función Gama: continuación análitica, ceros, polos y propiedades básicas.
%    \item La función Zeta de Riemann: representación integral, ecuación funcional y continuación analítica.
%    \item Ceros de la función Zeta, la Hipótesis de Riemann y su relación con el Teorema de los Números Primos.
%\end{enumerate}

\item \textbf{Tópicos adicionales (cubrir al menos uno a criterio de la persona docente) (2 semanas):} 
\begin{enumerate}
    \item El Teorema de Factorización de Hadamard.
    \item La función Gamma y la función Zeta de Riemann.
    \item Funciones elípticas.
    \item Los Teoremas de Picard.
    \item El Teorema de la Aplicación Conforme de Riemann.
    \item La Transformada de Fourier y aplicaciones.
\end{enumerate}

\end{enumerate}

\subsubsection*{Metodología}

\noindent \textbf{Lineamientos metodológicos:} La persona docente a cargo de este curso debe respetar las siguientes pautas:
\begin{enumerate}
    \item Fomentar el desarrollo de intuición sobre los conceptos estudiados en el curso.
    \item Promover la comunicación clara y rigurosa de argumentos e ideas matemáticas de manera oral y escrita.
    \item Cuando la naturaleza del contenido lo permita, se deben utilizar representaciones visuales que apoyen la comprensión de los temas. Por ejemplo se puede consultar de referencia el libro \cite{Nee23}.
    \item Fomentar una visión integral de las matemáticas y de cómo se enlazan las diferentes áreas a través del análisis complejo.
\end{enumerate}

\textbf{Sugerencias metodológicas:} Adicionalmente, se tienen las siguientes recomendaciones para la persona docente a cargo de este curso:
\begin{enumerate}
    \item Utilizar el recurso tecnológico para reforzar distintos conceptos estudiados en el curso por medio de la visualización, cálculo y experimentación.
    \item Aprovechar momentos adecuados del curso para motivar el estudio por medio de reseñas acerca del desarrollo histórico del análisis complejo y su importancia en aplicaciones dentro y fuera de la matemática. \textcolor{red}{Agregar algunas referencias para historia del análisis complejo.}
    \item En la evaluación, se sugiere utilizar estrategias que promuevan el trabajo constante de parte de las personas estudiantes a lo largo del ciclo lectivo. Por ejemplo, esto se puede hacer por medio de listas de ejercicios recomendados o proyectos.
    \item Para fomentar el desarrollo de las habilidades investigativas en el estudiantado, se sugiere la implementación de proyectos de investigación y participación en coloquios o seminarios de investigación.
    \item Dado que hoy en día la mayoría del trabajo en matemática, tanto a nivel académico como profesional, es de naturaleza colaborativa, se sugiere a la persona docente implementar estrategias que promuevan el trabajo en equipo.
\end{enumerate}


\nocite{AF03}
\nocite{AF21}
\nocite{Ahl21}
\nocite{Apo76}
\nocite{BN10}
\nocite{Boa10}
\nocite{Con78}
\nocite{Con95}
\nocite{GK06} 
\nocite{Lan99}
\nocite{Kra08}
\nocite{Nee23}
\nocite{Pri03}
\nocite{SS203}




\printbibliography[heading=subbibliography,section=\the\value{refsection}]
\end{refsection}
